\documentclass[12pt, a4paper]{book}%, oneside

\usepackage{times}
\usepackage[headings]{fullpage}
\usepackage{setspace}
\usepackage{listings}
\usepackage{color}

%\definecolor{mygreen}{rgb}{0,0.6,0}
%\definecolor{mygray}{rgb}{0.5,0.5,0.5}
%\definecolor{mymauve}{rgb}{0.58,0,0.82}

\lstdefinestyle{csharp}{basicstyle=\ttfamily\footnotesize,language=[Sharp]C,breaklines=true,commentstyle=\itshape,escapeinside={/@}{@/},frame=lines,keywordstyle=\bfseries,keywordstyle=[2]\bfseries\itshape,numbers=left,stepnumber=2,numberstyle=\tiny}

\newcommand{\csharp}[2][]{\mbox{\lstinline[basicstyle=\ttfamily,language=[Sharp]C,commentstyle=\itshape,escapeinside={/@}{@/},keywordstyle=\bfseries,keywordstyle=[2]\bfseries\itshape,morekeywords={[2]{#1}}]!#2!}}

%\lstset{ %
%  backgroundcolor=\color{white},   % choose the background color; you must add \usepackage{color} or \usepackage{xcolor}
%  basicstyle=\ttfamily\footnotesize,        % the size of the fonts that are used for the code
%  breakatwhitespace=false,         % sets if automatic breaks should only happen at whitespace
%  breaklines=true,                 % sets automatic line breaking
%  captionpos=b,                    % sets the caption-position to bottom
%  commentstyle=\color{mygreen},    % comment style
%  deletekeywords={...},            % if you want to delete keywords from the given language
%  escapeinside={\%*}{*)},          % if you want to add LaTeX within your code
%  extendedchars=true,              % lets you use non-ASCII characters; for 8-bits encodings only, does not work with UTF-8
%  frame=single,                    % adds a frame around the code
%  keepspaces=true,                 % keeps spaces in text, useful for keeping indentation of code (possibly needs columns=flexible)
%  keywordstyle=\color{blue},       % keyword style
%  language=[Sharp]C,                 % the language of the code
%  morekeywords={*,...},            % if you want to add more keywords to the set
%  numbers=left,                    % where to put the line-numbers; possible values are (none, left, right)
%  numbersep=5pt,                   % how far the line-numbers are from the code
%  numberstyle=\tiny\color{mygray}, % the style that is used for the line-numbers
%  rulecolor=\color{black},         % if not set, the frame-color may be changed on line-breaks within not-black text (e.g. comments (green here))
%  showspaces=false,                % show spaces everywhere adding particular underscores; it overrides 'showstringspaces'
%  showstringspaces=false,          % underline spaces within strings only
%  showtabs=false,                  % show tabs within strings adding particular underscores
%  stepnumber=2,                    % the step between two line-numbers. If it's 1, each line will be numbered
%  stringstyle=\color{mymauve},     % string literal style
%  tabsize=2,                       % sets default tabsize to 2 spaces
%  title=\lstname                   % show the filename of files included with \lstinputlisting; also try caption instead of title
%}

\usepackage{amsmath, amsthm, amssymb}
% \usepackage[top=1.5cm,right=1.5cm,bottom=1.5cm,left=1.5cm]{geometry}
% \usepackage{fancyhdr}
% \setlength{\headheight}{20pt} 
% \usepackage{graphics}
\usepackage{graphicx}
\usepackage{epstopdf}
\usepackage{float}
\usepackage{subfig}
\usepackage{array}
% \usepackage{xtab}
\usepackage{multirow}
\usepackage{booktabs}
\usepackage{appendix}
\usepackage[pdfborder={0 0 0}, colorlinks=true, linkcolor=black, citecolor=blue]{hyperref}
\usepackage[printonlyused]{acronym}
\usepackage{algorithmic}
\usepackage{algorithm}
\usepackage{ifpdf}
% \setlength{\parindent}{0pt}

\newcommand{\submissionDay}{04}
\newcommand{\submissionMonth}{June}
\newcommand{\submissionYear}{2013}
\newcommand{\submissionDate}{\submissionDay~\submissionMonth,~\submissionYear}
\newcommand{\typeOfThesis}{Bachelor Thesis}

\newcommand{\titleOfThesisOne}{Ju-Jutsu Training Kinect Application}

\newcommand{\authorOfThesis}{ElHassan Makled}
\newcommand{\supervisorOne}{Assoc. Prof. Georg Jung}

\newcommand{\includefig}[4]{
    \begin{figure}[ht]
     \centering
      \includegraphics[width=#1\textwidth]{images/#2}
      \caption{#3}
      \label{#4}
    \end{figure}
}

\newcommand{\includefigWSC}[5]{
    \begin{figure}[ht]
     \centering
      \includegraphics[width=#1\textwidth]{images/#2}
      \caption[#3]{#4}
      \label{#5}
    \end{figure}
}

\newcommand{\includeeps}[4]{
\includefig{#1}{#2.eps}{#3}{#4}
}

\newcommand{\includeepsWSC}[5]{
\includefigWSC{#1}{#2.eps}{#3}{#4}{#5}
}

\newcommand{\N} {\ensuremath {\vec{N}}}

\ifpdf
\pdfinfo {
	/Author (\authorOfThesis)
	/Title (\titleOfThesisOne)
	/Subject (\typeOfThesis)
	/Keywords ()
	/CreationDate (D:20090707085533)
}
\fi

\begin{document}
% \overfullrule=5pt
\pagestyle{plain}
\pagenumbering{Roman}

\newcommand{\titlePage}{

\thispagestyle{empty}
\begin{center}
	\textbf{Media Engineering and Technology Faculty}\\[1mm]
	\textbf{German University in Cairo}\\[1mm]
	\includegraphics[width=2.5cm]{GUC-logo-ss.eps}
	
	\vspace{2cm}
	\doublespacing
	{\Huge \textbf{\titleOfThesisOne}}\\
	\singlespacing
	\vspace{2cm}
	{\large \textbf{\typeOfThesis}}\\
	
	\vfill
	\parbox{1cm}{
  		\begin{large}
    			\begin{tabbing}
       			Author: \hspace{2cm}  
        			\=\authorOfThesis\\[2mm]
      			Supervisors: 
        			\>\supervisorOne\\[2mm]
				\>\supervisorTwo\\[2mm]
				\>\supervisorThree\\[2mm]
      			Submission Date: 
        			\>\submissionDate\\
    			\end{tabbing}
  		\end{large}
	}\\
\end{center}
\clearpage
}
%++++++++++++++++++++++++++++++++++++++++++++++++++++++++++++++++++++
\titlePage
\thispagestyle{empty}\ \clearpage
\titlePage
%++++++++++++++++++++++++++++++++++++++++++++++++++++++++++++++++++++
\thispagestyle{empty}
This is to certify that:
\begin{itemize}
\item[(i)] the thesis comprises only my original work toward the Bachelor Degree
\item[(ii)] due acknowlegement has been made in the text to all other material used
\end{itemize}

\vspace{2cm}
\begin{flushright}
\rule[0mm]{6cm}{0.2mm}\\
\authorOfThesis\\
\submissionDay~\submissionMonth,~\submissionYear\\
\end{flushright}
\clearpage


\chapter*{Acknowledgments}
\addcontentsline{toc}{chapter}{Acknowledgments}
\label{chap:ack}

This project would not have been possible without my supervisor, professor Georg Jung's continuous support and encouragement. I am also very grateful to my family because of the support they gave to me from the beginning of the project till the very end.\\

I would like to give a special thank you to my grandmother, who has motivated me to work when I needed it the most.\\

Lastly, I would like to thank my friends Omar Mowafi, Fady Kamal, Nada Nasr, Mai Nasser, Amr Tharwat, Tarek Samy and many others for giving me feedback about my work when I needed it, as well as supporting me the whole time.

\chapter*{Abstract}
% \addcontentsline{toc}{chapter}{Abstract}
\label{chap:abstract}
IMPACT is a multidisciplinary project in which students of computer science (more precisely: software engineering), material sciences, mechatronics and embedded system design come together to create Thai-boxing pads with builtin impact sensors and wireless connection to a base station.
In this project, we will implement a multi-sensor tracking and analysis tool for a physical workout routine specific to Thai/kick-boxing or similar contact sports to track and display performance and fitness level of a practitioner over a single and over multiple sessions. It will be used to record and correlate the inputs from impact sensors and technique recognition and track improvemets of practitioners over time.
\tableofcontents
% \addcontentsline{toc}{chapter}{Contents}
\cleardoublepage 

\pagestyle{headings}
\pagenumbering{arabic}

% \setlength\parskip{15pt}
\setlength\parskip{.5\baselineskip plus .2\baselineskip
	minus .4\baselineskip}
% \setlength\parskip{.5\baselineskip \@plus .1\baselineskip \@minus ..1\baselineskip}


\chapter{Introduction}
\label{chap:intro}
With the current rise of ubiquitous computing touchless input sensor devices have emerged to replace controllers in some video games, house hold appliances or lately, mobile phones. With these devices present, programmers take advantage of developing applications either for entertainment, such as dancing games, sword fighting games or shooting games, or fitness related applications, like yoga.  
\\
\\
When it comes to fitness related applications, users are able to use them as support for keeping track of their performance. Some applications target to replace a coach, such as Nike+ Kinect training. This program guides the user through a series of exercises including sit ups, push ups and more. The user in the practice is able to recieve real time feedback and coaching.
\\
\\
All the emerging fitness applications mostly target activities that could be done at a limited space such as a living room. Only few applications target contact sports such as Ju-Jutsu, Muay Thai, or Boxing. However, there applications available using touchless input sensors that target these types of sports, but only focuses on the gaming experience and do not act as fitness applications.
\\
\\
Ju-Jutsu, Muay Thai, Karate, Boxing, and other contact sports have a certain practice that requires a practitioner and a coach, where the coach is holding two Thai-Pads, a thick pad that covers the arms of the coach so that the trainee would punch or kick for practice. Usually this practice is fast paced with the coach always signalling for a certain move to be executed by the practitioner. With the practice being fast paced the practitioners sometimes find it hard to recognize their mistakes or evaluate their performance.
\\
\\  
Project Recon focuses on creating a tool for practitioners to evaluate and support in keeping track of their performance. Project Recon is part of a larger project titled Impact. Project Impact is a fitness monitor project that targets Ju-Jutsu practitioners. It has multiple plugins for applications. Those applications take input from sensors and process them. This information is later sent to Project Impact which will visualize the information to the user. Their are various types of information about the practitioner. All of which later reflect the performance and work out results of the user. The applications used as plug-ins to Project Impact include Project Recon. Which uses a Microsoft Kinect sensor to locate the practitioner and recognizes his techniques. Other possible sensors are the sensor equipped Thai-Pads and a heart rate monitor.

%\section{Section Name} \label{sec:s1}
%Some sample text with an \ac{ac}, some citation \cite{citeKey1}, and some more \ac{ac2}.

%\section{Another Section}
%Reference to Section \ref{sec:s1}, and reuse of \ac{ac} nad \ac{ac2} with also full use of \acf{ac2}.

\chapter{Background}\label{chap:background}
\section{Architecture}
The architectures shown below describes how communication between Project Recon and the general interface would look like. Figure 3.1 shows more details regarding the internal architecture and design of Project Recon itself.
\\
\\
As shown in figure 3.1, Project Recon main inputs are the Kinect sensor and Thai Pad Sensor. The Kinect sensor sends the Image stream and skeleton stream to Project Recon's Technique recognizer which in turn would take the skeleton of the user and keep track of its joints' position in time. When the practitioner moves the technique recognizer would compare the current movement with the database of reference moves. From this, it would associate the detected move with its respectful reference. The sensor equipped Thai Pads would measure the impact and send the value to Project Recon. The impact is then linked to its respectful technique and compared with reference impact. All this data is sent to the client which is responsible for the connection between the interface and Project Recon.
\\
\\
\centerline{\includegraphics[scale=0.4]{project_Recon_diagram.png}}
\\
\centerline{fig.3.1}
\\
\centerline{(Internal architecture and design of Project Recon)}
\\
\\
\\
In figure 3.2, an overview of the project is shown. Where it shows the linking between the Main Interface and Project Recon. As previously said, the client is responsible for the connection between both by communicating and connecting to a socket specified for it by the Main Interface.
\\
\\
\centerline{\includegraphics[scale=0.8]{general_overview.png}}
\\
\centerline{fig.3.2}
\\
\centerline{(Overall architecture of Project Recon and its connection to the interface)}

\section{Development Tools and Technologies}
When it came to developing with Kinect, two choices were on the table. Whether to use the official Microsoft Kinect SDK or use other unofficial SDK's. The official Kinect SDK may have slightly less capabilities, and we will be talking about them later and exploring them more. However, it is easier to use the official Microsoft SDK. The environment that was chosen for development (XNA) would also support our choice of using the official SDK. The programing language possible to use are Visual Basic and C\#. The chosen language was C\# as it was easier for me to develop for.

\subsection{Microsoft XNA}
XNA is an IDE (Integrated Development Environment) created by Microsoft to help game development for their systems to be easier. XNA is based on .NET Framework and it creates applications for platforms like XBox 360, Zun, Windows Phone and Microsoft Windows. Since its release there are 4 versions with the final version being XNA Game Studio 4.0. The framework was released on March of 2004 and the latest update released on September 16th 2010.
\\
\\
XNA is mostly used for XBox 360 video game development, which of course in turn is compatible with Kinect development, since Kinect is a Microsoft XBox 360 device. Making it easier to integrate the device within the code and collect data sent from the device and being able to read this data.
\\
\\
XNA makes it easier for developers to organize their code in many ways as it takes care of low level technologies related in game development. Helping developers to focus more on the gameplay and game detail itself. To make it more clear, when creating a class using XNA. Certain methods are created under a class titled \emph{Game1} this class inherits from \emph{Microsoft.XNA.Framework.Game}. This class includes several methods that help in organizing the code.
\\
\\
\emph{Initialize()}, in this method every single variable is initialized. In our case since we use the Kinect, several initialization took place. For example, initializing the variable kinect to a connected Kinect Sensor. As well as initializing the elevation angle of the Kinect itself. Every other variable is initialized in this method.
\\
\\
Code:
\begin{verbatim}
protected override void Initialize()
{
// TODO: Add your initialization logic here
kinect = KinectSensor.KinectSensors[0];
kinect.ColorStream.Enable(ColorImageFormat.RgbResolution640x480Fps30);
kinect.DepthStream.Enable(DepthImageFormat.Resolution320x240Fps30);
kinect.SkeletonStream.Enable();
kinect.Start();
kinect.ElevationAngle = 0;
colorData = new byte[640 * 480 * 4];
colorTex = new Texture2D(GraphicsDevice, 640, 480);
rawSkeletons = new Skeleton[kinect.SkeletonStream.FrameSkeletonArrayLength];
keepcover = true;
base.Initialize();
}
\end{verbatim}
The \emph{LoadContent()} method as its name specifies is responsible for loading the content required for the game. These contents range from images to audio or even video content that will be used in game as textures or in game music.
\\
\\
Code:
\begin{verbatim}
// Create a new SpriteBatch, which can be used to draw textures.
spriteBatch = new SpriteBatch(GraphicsDevice);
circleTex = Content.Load<Texture2D>("circle");
lineTex = Content.Load<Texture2D>("4KWjQ");
lineTexRed = Content.Load<Texture2D>("4KWjQR");
font = Content.Load<SpriteFont>("MainFont");
// TODO: use this.Content to load your game content here
\end{verbatim}
There is also an \emph{UnloadContent()} method that is used to unload the content once the game is done using them. This method is helpful with multi level games as loading all the game content unto the gpu will be impossible so items will be loaded when needed and unloaded when they are no longer needed during the gameplay to give space for more items to be loaded. In our project we will not need the Unload method as there aren't any high processing going on the graphics level.
\\
\\
The \emph{Update()} method takes the gameTime as an argument, which is a snapshot of the game timing state. This method is called every interval. Updating all variables and values that have changed from previous frames. In our case, the positions of the joints of the skeleton. Later, the draw method would take care in rendering the images and displaying them. The two figures below show two different frames in time with the skeleton in a position and in the other time in another position.
\\
\\
In the following images, two frames are shown where the skeleton first is standing regularly and then, later in time standing on one foot and leaning slightly to its left. As said before the update method updated the data from the previous method in terms of each joint's location in space and then the draw method rendered the skeleton with the new joints' locations.
\\
\\
\centerline{\includegraphics[scale=0.5]{skeleton_frame1.png}}
\\
\centerline{fig.4.1 (initial frame)}
\\
\centerline{\includegraphics[scale=0.5]{skeleton_frame2.png}}
\\
\centerline{fig.4.2 (few seconds later from initial frame)}
\\

\subsection{Kinect}
What makes Kinect sensor a good choice for this project is its ability to recognize and keep track of a person's movement through creating a virtual skeleton that matches the user and makes the skeleton move with the user. There are few constraints unfortunately when it comes to incorporating it with the Jiu Jitsu training. One of which would be the fact that Kinect does not differentiate between a user who is facing the Kinect from a user who is giving his back to the Kinect sensor.
\\
\\
The Kinect sensor consists of three major parts. A regular camera used mostly for video chat on the XBox and the Kinect and sometimes visualization (e.g adding the user inside the scene of the game). The other two major parts are the ones that detect the user and their motion. The first is an infrared emitter and the other is an infrared sensor. The way the Kinect works is that the infrared emitter emits a mesh of infrared rays that would reflect back to the Kinect off from the user and then are detected by the sensor.
\\
\\
The Kinect sensor has sends as an output three streams, an audio stream, a depth image stream and a color image stream. In the image below the color image stream is used to render the regular image captured by the Kinect camera.
\\
\centerline{\includegraphics[scale=0.5]{color_data.png}}
\\
\centerline{fig.4.3 (Color stream data visualization)}
\\
\\
Other streams that the Kinect provides as output are the skeleton stream, Kinect provides an array of skeletons (having skeletons of all available users). So far the array has a maximum of four users at a time. The skeleton stream can be used to acquire information regarding the skeleton of a certain user. In the next image, a circle was used to identify the places of the joints and was drawn on its respective joint. Different shades of grey were used in this image to show the system's differentiation of the joints.
\\
\centerline{\includegraphics[scale=0.5]{color_joints.png}}
\\
\centerline{fig.4.4 (Color stream data visualization with a visualization}
\\
\centerline{of the joints that is provided from the skeleton stream)}
\\
\\
There is a total of 20 joints recognized by the Kinect sensor. The head, the center shoulder, right shoulder, left shoulder, right and left elbows, right and left wrists, right and left hands, right and left hips, right and left knees, right and left ankles, right and left feet, a center hip and a spine.
\\
\\
After visualizing the joints, the next step was to link the joints with a line to show the skeleton itself. By simple line equations and using the two points as the joints, the skeleton was drawn in the following manner:
\\
\\
\centerline{\includegraphics[scale=0.5]{color_skeleton.png}}
\\
\centerline{fig.4.5 (Color stream data visualization and skeleton visualization)}
\\
\\
\subsection{Connections and Communications}
The communication between Project Recon and the interface is seen as a client-server communication. When browsing the possible technologies to use, socket programming came at the top of the list.
\\
\\
Socket Programming uses the client server model of network communication, in which a server (in this case the interface) would listen on different ports for data arriving from any of the devices that are supposed to connect to it. When Project Recon starts and is ready it will take the part of the client, as it will send its data stream to the specified socket. When ready, a handshake will occur to initialize connection between starting to send the data stream. In which the interface will receive the data of gestures and delivers the proper output.
\\
\\
There are three different data streams, a data stream responsible for system-related gestures, like a gesture to start the session or end it. The other stream consists of technique details. Finally, the third stream is contains the impact of the technique. All of the three streams have a timestamp in order to associate each impact with the proper technique used.

%Project Fast in Recognition in real time
%Robust and Dynamic in terms of Motion Recognition
%Be able to plugin to different systems
%be able to send data to the interface ad monitor in order to make it visibe to the user playing

\section{Project Requirements}
%In order to ensure best performance, certain constraints need to be handled. 
The practice that is targeted by the project is fast paced, which means the system needs to be performant enough to handle fast motion and recognize the technique in real time. Also, for it being a dynamic practice, which means rapid movement takes place during practice and the user will not necessiraly be always facing the Kinect as he will be moving freely around the practice room and most of the times will be giving his back to the Kinect. We need to make the system robust enough to be able to recognize different motions.
\\
\\
Another basic requirement is to allow the differentiation between two different users infront of the Kinect, one being the coach and the second being the practitioner. This is easily resolved as Kinect's stream sends multiple skeletons of the multiple users. However it is required that the first user be defined as a the practitioner and the second to join in later would be the coach.
\\
\\
Finally, the project should be able to plugin to a different system, the fitness monitor. The fitness monitor is an interface that shows the user's current stats, technique executed, and more. A connection should be established between Project Recon and the fitness monitor, where Project Recon should be able to send information regarding the technique executed by the user for the fitnes monitor display.
\chapter{Application}\label{chap:application}
In order to create a robust system, two solutions are proposed. The first utilizes the fact that a person faces 
the same direction their waist faces. So we create a plane between three fixed 
points, the center hip $\vec{c}$, the left $\vec{l}$ and the right hips $\vec{r}$. A normal vector \N was created 
from this plane in order to allow the program to recognize where the person is 
facing.

\begin{figure}[!htbp]
\centering
\includegraphics[width=1\textwidth]{images/skeleton_frame3.png}
\caption{In the figure the three joints are pointed out and named}
\label{skeletonframe3}
\end{figure}

To calculate the normal vector, we subtract the right and left hip joints from 
the center hip joint and then we calculate the cross product of the two 
subtractions.

$\vec{N}$ = ($\vec{r}$ - $\vec{c}$) $\times$ ($\vec{l}$ - $\vec{c}$)

We then calculate the normal by using the Normalize method.

\begin{verbatim}
var dir = Vector3.Cross(rhip - chip, lhip - chip);
var norm = Vector3.Normalize(dir);
\end{verbatim}

As the person turns, the normal vector will be used as a reference to calculate the angle between it and the Z axis. This will allow us to rotate the skeleton so that it will be always facing the Kinect. Recognition will then take place.

Kinect has another problem, which is that it cannot diffrentiate between a user giving it his face or his back. In other words, a user's left and right are swaped if the user is giving the Kinect his back.

%put an image to make it better understandable

The normal vector calculation approach will reswap the left joints with the right joints when the \N is on the negative Z axis. 
However, here another problem emerges. When the swap takes place, an error occurs. Mainly, the programmer can not 
edit the value of the position of the joints. The value is read from the skeleton 
stream of the Kinect. In order to resolve this, an avatar skeleton needs to be 
created, littleGuy. littleGuy is an array of SkeletonPoints, where every 
SkeletonPoint represents a joint, and gets the value from these joints. This 
helps in allowing us to edit and the SkeletonPoints as we wish and in the end 
draw this avatar skeleton instead.

In figure~\ref{skeletonnormal} the user is facing their right, having the normal vector (red 
line) erected and defining the direction the user is facing.

\begin{figure}[!htbp]
\centering
\includegraphics[width=1\textwidth]{images/skeleton_normal.png}
\caption{Skeleton with a normal vector defining where it is facing}
\label{skeletonnormal}
\end{figure}

The translation and rotation of the skeleton occurs in the Skeleton Space. The algorithm translates littleGuy to the origin point where it calculates the \N. Rotation occurs in the origin and then the new manipulated littleGuy translates back to a fixed point in the Z axis, the Y axis and X axis however remain in the origin. After this Project Recon stores the translated littleGuy in a new avatar skeleton called transGuy. The program then renders transGuy to the 2D Color Space. 

The second step is using Kinect's face detection capabilities, which may be 
hard and require high processing. Errors may occur when Kinect re-swaps the joints to their right place. For this we will use the face dectetion capabilities of Kinect to create a sanity check. That is if the swap did occur and the user's face was visible, then the swap is a mistake and it will re-swap again to the correct state.

When it came to testing the normalization of the skeleton using the Kinect, a problem emerges. When the user has a joint not facing the Kinect, then the Kinect does not have the position value of this joint. The reason for this is that no infrared rays are reflected from this joint. In turn, when the rotation occurs the some joints have false space(X, Y, and Z) values. This error will question the validity of the rendered and final translated values of the joints. This error made us discard this method and resort to a different one that will also aims to maintain the robustness of the system.

%here talk about the final way and discuss how you will use the recognition to solve the issue with robustness 
\section{Recognition}
Recognition is the essence of this project. Recognition needs to be in real time, but most importantly how will the recognition occur. We discussed in the previous section the resolve of the normal vector. However, why is it needed? In recognition we want to be able to translate the skeleton of the person in order to always recognize from the same direction, facing the Kinect. In the coming sections we will discuss the different ways for recognition.

\subsection{Glyphs Method}
This is the proposed method in the case of recognition taking place after rotation by the help of \N as discussed earlier in this chapter. The idea behind it is to create a path taken by each joint from the skeleton and store it in an image. Creating a single image with the overall motion is and comparing it with the reference image, would be more optimal than taking a number of frames and comparing them with other reference frames.

As a joint moves, it creates a path shadowing its past movements. Each joint has its own exclusive color in order to differentiate if lines intersect. When the user moves the system will draw the paths taken by the joints creating some sort of a glyph image depicting the overall motion of the skeleton. An example of how a glyph would like is in both figures below. Figure 1 shows the user as he executes a move with each joint creating a path. Figure 2 shows the final glyph as it would be stored.

%image of a person with a glyph
%image with a glyph

After storing the glyph image, a process of image processing would take place where comparisons between the captured glyph image and the reference glyph images would begin. When a match is found, the technique would be recognized, if a match was not found the glyph image is discarded. Recognition when comparing two glyph images would of course have some sort of threshold as users are different in their executions of a technique.

Limitations that would emerge from using this method are the fact that motions on the z axis would be hard to recognize as the path would on the z axis wont be drawn. To solve this limitation, it is recommended that the rotation of \N around the y axis change. Rather than rotating the \N to the z axis it would rotate to form a 45$^\circ$ on the xz plane of the positive z and x axis. This way the glyph image would have a representation to motion detected on all three axes.

It should be clear as to when the program should record the path patterns and store the glyph. In order to know when exactly to record, we should first understand when will techniques be executed. As stated earlier, the coach gives signals to the practitioner for him to execute specific moves. Since the coach will be considered a user, it is recommended that recording of path paterns takes place as soon as a signal is initiated by the coach.

Since we will be unable to use the method of normalizing the skeleton due to the limitations of the Kinect sensor, this method would be obsolete in our application as the user is mobil and active around the room. Making it hard to have standard glyph patterns.

\subsection{Joint Position Lists}
This method focuses on creating a list of the last 30 positions of a joint. Where the system stores the newest position to an already full list, the last position(the earliest element on the list) is dicarded and replaced by the newest, see figure %put figure here describing the stack. 
Since the application's frame-rate is 30 frames per second. The system takes every 30 positions each second and between each two frames it has inbetween lists, see figure  %put figure here describing the 30 and 30

We create a class, StoreGesture, that describes the element of the list. StoreGesture has two attributes. Position, which is a 3D vector describing the joint's position. The second is Time, which is the specific time on which that specific position takes place. The code~\ref{lst:storegesture}shows the class StoreGesture.

\lstinputlisting[style=csharp, caption ={StoreGesture Class},label={lst:storegesture},float={tbp}]{listings/StoreGesture.cs}

A new class  DetectGesture is responsible for creating the list of StoreGesture and detecting the gestures. The way this works is that a list, posList is set to contain a list of StoreGesture and a maximum size of the list is set, which always is 30. In DetectGesture there are methods that would get the values inside the list, add an element to the list, and use the list to recognize a gesture.

Each different gesture has a series of joints that accent this gesture. Meaning that if the gesture is a punch, then the joints that are important and mostly resemble the punch are the elbow and wrist joints. Same thing for a kick, the knee's and the ankle's positions are important in order to recognize the gesture.

As shown in listing~\ref{lst:jointsposlist}, for each joint that would be later detected we create a DetectGesture with a size of 30 frames in the main class of Project Recon.

\lstinputlisting[style=csharp, firstline = 26, lastline = 36, caption ={Joints Position Lists},label={lst:jointsposlist},float={tbp}]{listings/xna.cs}

As the code updates, we update this list by adding the new position of the joint to its respective list by using the method in DetectGesture known as addPosition. This method takes as arguments a SkeletonPoint and a Kinect sensor. It creates a new StoreGesture entry with the given SkeletonPoint in its arguments. Then, it checks if list is full. If the list is full it would remove the earliest entery. Then in the end it would add to the list the new entry. Listing~\ref{lst:addposition} shows the addPosition method.

\lstinputlisting[style=csharp, firstline = 1, lastline = 13, caption ={addPosition Method},label={lst:addposition},float={tbp}]{listings/DetectGesture.cs}

As discussed earlier, the system needs to be robust enough to recognize gestures done at any angle the user is facing from or away of the Kinect. In order to acheive this, gesture detection needs to be exclude relations with the z and x axis as they will be dynamically changing with the time. Instead, a good way to tackle this is to detect a gesture through analyzing relations between joints and each other. For example, in a specific kick (known as the forward thrust) the knee will be higher than the hips in most cases. So taking advantage of this relation will help in acheiving the robustness of the system that is required.

The Detect method in DetectGesutre, takes two arguments, reference1 and reference 2. Those arguments resemble two different position lists of reference joints that will be used in order to aquire relationship between them and the main joint. In order to recognize a gesture, we take divide a gesture into three different parts. Let us take a punch for example. In figure %punch figure
we divide the punch into three main motion frames. By observing the frames one can simply find a relation between the joints and each other. In this example the elbow is closer to the spine in the first and third frame in terms of the y axis, however closer to the shoulder in the second frame. Overall, another added information is that the angle between the elbow-shoulder and shoulder spine in a punch is more than 45$^\circ$. 

Similarly, the Detect method divides the frames into 3 different parts, each with 10 elements in them. A loop goes through them and whenever one element fits the criteria of one of the parts, a value that is responsible for calculating the probability is incremented. After the loop goes through the whole 30 frames. The probability value is checked, where 30 is 100 percent match. A threshold is given that is 25. Depending on the probability value either passing or not passing the threshold the gesture will be detected or not detected. Listing~\ref{lst:detect} shows the Detect method.

\lstinputlisting[style=csharp, firstline = 15, lastline = 74, caption ={Detect Method},label={lst:detect},float={tbp}]{listings/xna.cs}

Since every $\frac{1}{3}$rd of a second the class adds a position and removes another and then checks the latest list with the new position, then a gesture might be recognized more than once. In order to avoid this, we define a minimal period between gestures. The list becomes a completely new list every second. So the minimal period between gestures is 1 second. However, what if the person is fast enough to execute 2 gestures in less than two seconds. If this is the case and the system recognizes a gesture before the minimal period between gestures elapse. The system will compare it with the last gesture recognized. If they are the same, then it is discarded and not taken as a gesture. If not then the system considers it a gesture.

\subsection{MCS UK Solution Development Gesture Service \cite{citeKey3}}
%DISCUSS HERE Microsoft Consulting Services gesture service and if it would be used how to acheive robustness
Microsoft Consulting Services UK (MCS UK) have created a gesture service that Kinect SDK users who want to create an application that uses gestures could adopt. The gesture service is written in C\#. It analyzes the gestures and recognizes them.

It is very similar to joint position lists in terms that it defines a gesture by frames or parts. As figure~\ref{wave} shows two parts of a gesture.

\begin{figure}[!htbp]
\centering
\includegraphics[width=1\textwidth]{images/wave.png}
\caption{An image from MCS UK's blog describing two parts of a wave}
\label{wave}
\end{figure}
 
They further describe this as not enough in order to accurately recognize a gesture. Stating that if the user who is waving drops his hand between the two parts, then the system would still recognize the gesture even though the actual gesture was not executed.

To resolve this problem they view the different gesture parts as different states. In figure~\ref{gesturestates} it shows three gesture parts, each resembling a state. Currently the user is in gesture part 2. There are three results, either fail, pause or succeed. Depending on certain situations these results influence the next state the system will go to.

\begin{figure}[!htbp]
\centering
\includegraphics[width=1\textwidth]{images/gesturestate.png}
\caption{Three gesture parts resembling states}
\label{gesturestates}
\end{figure}

The fail result occurs when the user moves in an inconsistent way with the full gesture. This result will return to the begining of the gesture and start from the first gesture part. The Pausing result occurs when the user moves too slow or they did not perform the next part of the gesture. The system pauses for a maximum of 100 times. If the gesture remains and the next state was not executed then the gesture automatically fails and recognition returns from the first state again. Succeed takes place when the user executes the current part of the gesture. After a short pause the system would wait for an execution of the next gesture part from the user.

\subsubsection{Architecture of the MCS Gesture Service}
The gesture service consists of three parts. Gesture Controller, Gesture, and Gesture Parts. Figure~\ref{gestureservicearc} visualizes the architecture.

\begin{figure}[!htbp]
\centering
\includegraphics{images/mcsarc.png}
\caption{An image from MCS UK's blog describing the architecture of the service}
\label{gestureservicearc}
\end{figure}

The gesture controller is responsible for controlling all the gestures available for the user to perform. The gesture is responsible for controlling the gesture parts as well as keeping track of the gesture that is considered in the current state. The gesture block contains an array of IRelativeGestureSegment which are individual implementations from the IRelativeGestureSegment interface. A skeleton frame is created and checked through the states. Once the final state returs `Succeed' it raises a gesture recognized event. This event is caught by the gesture controller. The IRelativeGestureSegments resemble the individual parts of the gesture.

To handle robustness when using this method. Similar to the previous method, we take joint positions relevent to other joint points and not the axes.
\section{Connecting to the Interface}
After recognition takes place and Project Recon recognizes a gesture. Project Recon needs to send the gesture and the time frame which the gesture took place to the interface so that the interface would view it to the user in real time during practice and also keep track of his performance.

Two approaches were considered. The first focuses on using socket programming to create a server and client side where communication would occur on certain ports. The second was to use a service known as Pusher.
\subsection{Pusher \cite{citeKey4}}
Pusher is an API for creating quick realtime connection between twointernet connected devices through WebSockets.

Pusher provides many libraries for differerent frameworks like JavaScript for HTML apps or .NET and python. To better understand Pusher, figure~\ref{pusher} shows a flow diagram of Pusher.

\begin{figure}[!htbp]
\centering
\includegraphics{images/pusher.png}
\caption{An image from Pusher's official website describing the data flow using the API}
\label{pusher}
\end{figure}

Pusher works as a Publish/Subscribe model. Where users create channels and applications connect to this channel to initiate a bi-directional connection between them.

The problem with Pusher was that it would require the user to acquire internet access at all time. Since this will not always be the case and users may be offline we will use regular socket programming.
\subsection{Socket Programming}
This approach is simple and based on the server and client connection. Where the client side will be Project Recon. This approach also helps in creating an offline connection between both applications, unlike Pusher which requires internet access on the system. From Project Recon's side we create a class, client, that will function as the link between Project Recon itself and the interface.

The class, as shown in listing~\ref{lst:client}, has four methods. One method to establish the connection, connect. The second to send data to the server, send. Finally, the method close which is responsible for closing the connection between the server and the client.

\lstinputlisting[style=csharp, caption ={Client Class},label={lst:client},float={tbp}]{listings/client.cs}

The connect method creates a new client socket and initializes a port that will be used to send the data stream to. It then aquires the IP address and creates a remote end point where to send to. Since both applications (Project Recon and the Interface), then they will both have the same IP address. After defining the remote end point we connect the socket to it.

Now that the connection is established, we send the first confirmation data. Which is a string confirming that connection is established and the Kinect is ready for inputs from the user. Before sending the data, we encode the data to bytes and call the method send and give it the new encoded data.

The method send takes three arguments two of which are the sender and the data. The method encodes the data to byte and then sends the data.

After the Kinect Sensor is connected to Project Recon, Project Recon will directly attempt to connect to the interface. Once connection is established, Project Recon will run regularly. Once a gesture or technique is recognized Project Recon will use the client method send to send a string of a certain format. The format is as shown in figure~\ref{format}.

\begin{figure}[!htb]
\centering
\includegraphics[width=1\textwidth]{images/format.png}
\caption{Format of the data stream sent to the interface}
\label{format}
\end{figure}

The interface will take the byte data stream and decod it to a string. It will take the time of the gesture or motion execution and from the data the interface will take certain actions.
% add more chapters here

\chapter{Conclusion}\label{chap:concl}
%Project Fast in Recognition in real time
%Robust and Dynamic in terms of Motion Recognition
%Be able to plugin to different systems
%be able to send data to the interface ad monitor in order to make it visibe to the user playing
The purpose of Project Recon as stated earlier is to create an application that recognizes Ju-Jutsu moves during a practice session in real time. The application has to give feedback on technique executions in terms of body movement, speed, and impact. To fulfill this purpose the application requires to be robust and dynamic in recognizing techniques. Another requirement was to recognize multiple users, a coach and trainee. Finally, the project should be able to plugin to the interface and be able to establish a bi-directional communication stream between it and the interface.

Using the Kinect sensor, XNA, and Kinect SDK we were able to create a robust system that will be able to recognize a technique executed. Also we were able to establish a connection using Socket programming to define a connection between the two applications, Project Recon and the interface.

Due to various limitations of Kinect, the complete project requirements were not reached. The scope of the project was beyond the capabilities of the Kinect.

The complete robustness of the project when it comes to recognition is not acheived. As well as the accuracy regarding the Kinect to recognize minor motion that could determine the technique quality of execution. Hence, the feedback that is required from the project to deliver is not achieved.

To better understand the limitation, let us take a punching technique as an example. To execute a proper punch, the trainee twists his hip to the direction of the punch as his arm gives the blow. The trainee's palm begins by being close to the torso facing upwards and ends by being completely stretched and facing downards. The reason for this limitation is that Kinect will not be sensitive enough to recognize these minor moves of the palm facing downwards or upwards or the hips twisting with the punch. Another reason is that Kinect recognize only 20 joints of the human body, which makes it hard for it to acquire accurate values of a specific motion.
\chapter{Future Work}
\label{chap:todo}

In this section, I will list some ideas and expansions to the tools that could be implemented in future versions of the impact project.

\section{Trainers integration}
In this feature, trainers will have their own accounts where they can handle multiple trainees and assign some sessions for them. They can monitor the level of their trainees, as well as giving feedback regarding their performance.

\section{Social training platform}
The idea of this expansion is that the tool could be an online social platform where practitioners share their sessions and experience with other fellow practitioners.

\section{Gamification}
This will add a competition-like environment to the trainees by having levels (amateur, expert,etc..), competitions between them and other features.

\newpage

\section{Plugins}
There are many great ideas for plugins that will help users.

\paragraph{Heartbeat rate sensor}
The heartbeat rate sensor could be used during the round or off-session to measure the heartbeat rate of the trainees to give an image of their fitness level.

\paragraph{Leap Motion sensor} \footnote{Leap motion main site https://www.leapmotion.com/}
This sensor is a sensor that looks like the xbox Kinect sensor. It will help get more accurate view of the position of the trainee during the session.

\nocite{*}
\appendix
\renewcommand{\appendixtocname}{Appendix}
\renewcommand{\appendixpagename}{\appendixtocname}
\addappheadtotoc
\setboolean{@twoside}{false}
\appendixpage

\chapter{Lists}
\addcontentsline{toc}{section}{List of Abbreviations}
\begin{acronym}[\hspace{3cm}]
  \acro{ac}[AC]{Acronym Without Citation}
  \acro{ac2}[AC2]{Acronym With Citation \cite{citeKey2}}
\end{acronym}
\clearpage
\listoffigures
\addcontentsline{toc}{section}{List of Figures}
% \listoftables
% \addcontentsline{toc}{section}{List of Tables}


\bibliographystyle{plain}
\bibliography{bachelor}
\addcontentsline{toc}{chapter}{References}

\end{document}
